% HSLU Thesis Format per WritingDissertations_21M.pdf (APA 7, British English)
% Lucerne University of Applied Sciences and Arts

\documentclass[10pt,a4paper]{report}
\usepackage[utf8]{inputenc}
\usepackage[T1]{fontenc}
\usepackage{mathptmx}
% Minimal margins to save space (HSLU minimum: left 4cm for binding)
\usepackage[a4paper, left=3.2cm, right=1.8cm, top=1.8cm, bottom=1.8cm]{geometry}

\usepackage{setspace}
\linespread{1.15} % Minimal line spacing to save space

\usepackage{booktabs}
\usepackage{tabularx}
\usepackage{graphicx}
\usepackage{amsmath}
\usepackage{amssymb}
\usepackage{hyperref}
\usepackage{natbib}
\usepackage{caption}
\usepackage{longtable}
\usepackage{array}
\usepackage{float}
\usepackage[most]{tcolorbox}
\usepackage{enumitem}

% Unicode/math support
\usepackage{newunicodechar}
\newunicodechar{→}{$\rightarrow$}
\newunicodechar{≥}{$\geq$}
\newunicodechar{≤}{$\leq$}
\newunicodechar{⁻}{$^{-}$}

% Graphics path: paper/ is current dir; figures copied here and also found in big5loop/
\graphicspath{{./}{../big5loop/visualization/figures/}{../big5loop/evaluation_data/personage/processed/}{../big5loop/evaluation_data/processed/}}

% APA-style headings per HSLU guide
\usepackage{titlesec}
\titleformat{\chapter}[display]{\centering\normalfont\bfseries}{\chaptertitlename\ \thechapter}{0pt}{\Large}
\titleformat{\section}{\normalfont\bfseries}{\thesection}{1em}{}
\titleformat{\subsection}{\normalfont\bfseries}{\thesubsection}{1em}{}
\titleformat{\subsubsection}[runin]{\normalfont\bfseries}{\thesubsubsection}{1em}{}

% Running header: first 2-3 words of section (APA recommendation)
\usepackage{fancyhdr}
\pagestyle{fancy}
\fancyhf{}
\fancyhead[C]{\textit{\nouppercase{\leftmark}}}
\fancyfoot[C]{\thepage}
\renewcommand{\headrulewidth}{0.4pt}

% Keywords format: italic, lowercase
\newcommand{\keywords}[1]{\par\vspace{0.5em}\noindent\textit{Keywords:} #1}

\hypersetup{colorlinks=true, linkcolor=blue, citecolor=blue, urlcolor=blue}

\begin{document}

% ========== TITLE PAGE ==========
% Per HSLU §1.2: repeat cover + email, supervisor(s), advisor(s)
\begin{titlepage}
\centering
\vspace*{2cm}
{\LARGE\bfseries Adaptive LLM-Based Chatbot with Personality-Aware Dialogue for Human-Centered Applications: Design, Implementation, and Pilot Evaluation with Swiss Informal Home Caregiver Coaching\par}
\vspace{2cm}
{\Large Jiahua Duojie\par}
\vspace{0.5cm}
{\normalsize duojie.jiahua@stud.hslu.ch\par}
\vspace{1.5cm}
{\large MSc Applied Information and Data Science\par}
{\large Lucerne University of Applied Sciences and Arts (HSLU)\par}
\vspace{0.5cm}
\vspace{1.5cm}
{\normalsize Supervisor: Prof.\ Dr.\ Guang Lu \quad guang.lu@hslu.ch\par}
{\normalsize Co-Supervisor: Samuel Devdas \quad samuel.devdas@hslu.ch\par}
\vfill
{\large March 2026\par}
\end{titlepage}

% ========== ABSTRACT ==========
% Per HSLU §1.2, §6: brief comprehensive summary; Keywords italic, lowercase, comma-separated
\chapter*{Abstract}
\addcontentsline{toc}{chapter}{Abstract}

Digital coaching can adapt to what a user says---but not yet to \emph{who they are}. Current systems personalise through mood detection or conversation memory, leaving the stable personality traits that determine coaching receptivity---how someone tolerates uncertainty, responds to structure, or processes emotional stress---entirely unmodelled. This gap is consequential: high-Neuroticism users disengage from emotionally incongruent systems; Conscientiousness predicts structured-goal adherence; personality-coaching fit shapes whether an intervention is used at all.

This paper presents \textbf{Big5Loop}, a personality-aware coaching architecture that closes this gap by adding a third personalisation layer---continuous Big Five (OCEAN) trait inference---on top of mood tracking and conversation memory. Per-turn personality detection uses a calibrated few-shot Llama 3 70B Instruct detector \citep{Dubey2024} (\texttt{meta/llama3-70b-instruct} via NVIDIA API); a two-timescale exponential moving average (EMA) pipeline separates volatile session states from stable cross-session trait estimates, mirroring the state/trait distinction established by experience sampling methodology. Stabilised trait estimates are translated into interpretable behavioural directives that continuously modulate emotional tone, relevance framing, and responsiveness across four coaching modes (emotional support, practical education, policy navigation, mixed). Three formal pipeline invariants enforce that personality detection runs on every turn, evidence retrieval activates conditionally on factual need, and every policy claim reaches the user citation-grounded---simultaneously resolving the tension between adaptive style and factual safety. A five-layer hybrid intent router combines LLM classification, confidence gating, keyword fallback, and non-overridable safety override to handle the full spectrum from acute emotional distress to Swiss social insurance navigation.

External validity of the personality detector was established against the full PERSONAGE corpus ($n = 318$, human-annotated ground-truth Big Five ratings): statistically significant Pearson correlations for all five traits (Openness $r = 0.41$, Conscientiousness $r = 0.43$, Extraversion $r = 0.39$, Agreeableness $r = 0.60$, Neuroticism $r = 0.53$; mean $r = 0.47$, all $p < 10^{-12}$). A style-aware ensemble variant achieves mean $r = 0.60$ on a balanced subset, approaching human inter-rater reliability for personality judgements from brief text \citep{Mairesse2011}. Zero-shot prompting produced inverted Conscientiousness ($r = -0.40$), confirming that style-anchored few-shot calibration---not raw LLM capability---drives detection quality. All five traits show consistent positive direction; the detector never inverts a trait's polarity.

Evaluations of coaching quality (RQ2), grounding safety (RQ3), routing accuracy (RQ4), and system reliability (RQ5) are fully pre-specified with target thresholds and are pending execution against the deployed system. These validation streams---automated pillar test matrices, LLM-as-judge coaching quality scoring, and load benchmarks---are described in Chapter~4 and constitute the next phase of the research programme.

The system is fully implemented and deployed: a three-tier architecture couples a conversational frontend, a partially parallelised orchestration pipeline (N8N v2) with detection caching, and a vector-enabled persistence layer with dual audit logging. Design coherence is enforced through shared typed contracts between all modules, ensuring that personality state, retrieval evidence, and grounding outcomes are consistently propagated and auditable at every turn. Big5Loop is, to our knowledge, the first system to unify personality-modulated coaching and citation-grounded policy guidance within a single coherent, auditable pipeline. Swiss informal caregivers serve as the pilot use case: 700,000 individuals managing complex care responsibilities with no personalised digital support. The architecture is domain-agnostic and designed for direct reuse wherever coaching style and factual accuracy must coexist.

\keywords{personality-aware coaching, Big Five OCEAN traits, large language models, retrieval-augmented generation, exponential moving average, informal caregiving, adaptive dialogue systems, Swiss social policy}

% ========== TABLE OF CONTENTS ==========
% Per HSLU §1.2: Level 1 and 2 headings only
\tableofcontents

% ========== MAIN BODY ==========
\clearpage
\setcounter{page}{1}

\chapter{Introduction}
\label{ch:intro}

Effective coaching adapts to the individual---not merely in content but in style, pacing, and emotional register. Human coaches implicitly learn how a client processes information, tolerates uncertainty, handles stress, and prefers to receive feedback. This stable understanding of who the person is, rather than what the person said last session, is the basis of coaching effectiveness \citep{Grant2003, Stober2006}. As coaching moves into digital systems, the central design question is: what mechanisms can replicate this adaptation?

Current digital coaching applications personalize through three progressively deeper mechanisms. The first is \emph{transient state tracking}: systems such as Woebot \citep{Fitzpatrick2017} and Wysa \citep{Inkster2018} detect the user's current mood and adapt the immediate response---selecting a breathing exercise for acute anxiety or a cognitive reframing prompt for rumination. This is reactive by design; it addresses the user's state at a single moment without modelling the stable traits that determine how they prefer to receive support. The second is \emph{conversation memory}: LLM-based systems (ChatGPT Memory, mem0, LangChain Memory) store prior utterances and retrieve relevant facts across sessions. A system may recall that a caregiver mentioned a dementia diagnosis or asked about benefit eligibility. Memory enables continuity, but personalises \emph{content}---what to reference---rather than \emph{style}---how to communicate. It stores the user's history without modelling who they are. The third layer, largely absent from deployed systems, is \emph{stable individual difference modelling}: inferring personality traits as enduring characteristics that should shape coaching behaviour across all interactions.

This gap is consequential. Coaching effectiveness is not personality-neutral: Big Five traits predict coaching outcomes, with Openness predicting benefit from exploratory approaches, Conscientiousness from structured goal-setting, and Neuroticism from emotional scaffolding \citep{deHaan2013, Jones2016}. A meta-analysis of 24 longitudinal multisource feedback studies confirmed that individual-level factors---including feedback orientation and perceived need to change---moderate whether and how much performance improvement follows structured feedback interventions, establishing personality as a boundary condition on intervention effectiveness that extends to coaching contexts \citep{Smither2005}. The same content delivered in different styles produces different outcomes depending on who receives it.

In informal caregiving specifically, personality shapes both how stress is experienced and which support strategies are effective. Neuroticism is a robust risk factor for elevated caregiver burden, depressive symptoms, and burnout: high-Neuroticism caregivers interpret demands as threats, exhibit poorer emotion regulation, and report greater subjective burden---a vulnerability that deepens over prolonged caregiving \citep{Penley2002, Brodaty2009}. Conscientiousness functions as a protective factor through structured problem-solving, adherence to self-care routines, and lower perceived exhaustion. Agreeableness facilitates adaptive coping via social support-seeking and collaborative provider interaction. These trait-driven differences are operationally significant: a high-Neuroticism caregiver experiencing benefit-application anxiety requires emotional validation before procedural guidance; a high-Conscientiousness caregiver in the same situation prefers a structured checklist. Personality-coaching fit is not a stylistic preference---it determines whether an intervention is used at all.

The Big Five framework (OCEAN: Openness, Conscientiousness, Extraversion, Agreeableness, Neuroticism) provides the validated theoretical basis for this mechanism \citep{Costa1992, McCrae1992}. Personality traits can be inferred from natural language with moderate to high accuracy \citep{Yarkoni2010, Park2015}, and personality-conditioned language generation improves perceived naturalness and satisfaction \citep{Mairesse2010, Devdas2025}.\footnote{\citet{Devdas2025} is an unpublished Master's thesis; the reported findings are available from the author on request.} The Zurich Model of Social Motivation \citep{Quirin2022, Quirin2023} provides a theoretical bridge from OCEAN traits to coaching behaviour (elaborated in \S\ref{ch:related}). The contribution of this work, however, is not personality detection but personality \emph{modulation}: translating stabilised trait estimates into behavioural directives that continuously adapt coaching across three quality dimensions---emotional tone appropriateness, relevance and coherence, and responsiveness to emotional needs.

Modulation must coexist with factual accuracy, a requirement current coaching systems sidestep by avoiding domain-specific knowledge delivery. In healthcare policy contexts, a coaching response must simultaneously validate a user's anxiety and provide citation-grounded procedural guidance. Rule-based CBT chatbots (Woebot, Wysa) deliver therapeutic exercises without domain retrieval. RAG-enhanced systems ground responses in evidence but apply no personality-aware modulation to how that evidence is presented. Ensuring that personality adaptation governs \emph{how} content is delivered without altering \emph{what} it says requires explicit architectural enforcement, not implicit reliance on the language model to maintain the distinction.

Swiss informal caregivers provide the concrete use case. Switzerland's Federal Statistical Office (FSO) reports life expectancy of 81.7 years for men and 85.7 for women (2022), with individuals aged 65+ already constituting nearly 20\% of the population---a proportion projected to exceed 25\% by 2040 as the ``oldest-old'' cohort (80+) is set to double by 2050 \citep{FSO2023}. The ratio of working-age adults to retirees has declined from 4:1 in 1990 to an anticipated 2:1 by 2040, creating a structural ``care gap'' in which informal caregivers currently provide approximately 80\% of all long-term care. Approximately 700,000 individuals provide this unpaid care to family members with chronic illness, disability, or age-related needs \citep{OECD2023, ZHAW2022}. They face four intersecting stressor domains: physical demands from direct care tasks and disrupted sleep; emotional burden from witnessing decline and managing challenging behaviours (e.g., dementia-related agitation, which affects 40--80\% of care recipients); financial strain from reduced employment and increased care costs; and social isolation from constrained participation and strained family relationships. Compounding this, Switzerland's Federal Health Insurance Act (KVG) covers only medically certified services and provides no direct compensation for informal caregiving; support is instead delegated to 26 cantons, creating a fragmented landscape where a caregiver in Zurich may access counselling and respite services while one in Appenzell receives virtually none. Swiss caregivers must navigate IV/Invalidenversicherung, AHV, Ergänzungsleistungen, cantonal Spitex services, and Hilflosenentschädigung in this heterogeneous policy environment. This domain simultaneously exercises all three personalisation layers---transient stress states, ongoing care memory, and stable personality traits---while demanding factually grounded policy guidance. It is thus an ideal proving ground for the present architecture. The system is designed for transfer to other coaching domains; Swiss caregivers are the use case, not the boundary.

Table~\ref{tab:gap} summarises the three-layer personalisation gap and positions Big5Loop against existing systems.

\begin{table}[H]
\centering\small
\caption{Personalisation layer gap summary}
\label{tab:gap}
\begin{tabularx}{\textwidth}{@{}lXlX@{}}
\toprule
Layer & Mechanism & Systems & Status \\
\midrule
Transient state & Mood/sentiment detection & Woebot, Wysa & Addressed \\
Content recall & Conversation memory & ChatGPT Memory, mem0 & Addressed \\
Stable individual differences & Personality trait modelling & None operational & Unaddressed \\
Coaching style adaptation & Trait-to-directive mapping & None operational & Unaddressed \\
Domain knowledge grounding & RAG with verification & Some medical systems & Not integrated with personality \\
Content/style separation & Architectural enforcement & None & Unaddressed \\
\bottomrule
\end{tabularx}
\end{table}

Mood tracking and memory are necessary---Big5Loop employs both---but insufficient for personalised coaching. No existing system combines personality-trait modulation, domain knowledge retrieval, and mandatory grounding verification in a pipeline with formal invariants.

This study designs, implements, and evaluates Big5Loop: a personality-aware coaching system that infers Big Five traits from multi-turn dialogue, stabilises them using EMA smoothing, and modulates coaching behaviour to improve emotional tone appropriateness, relevance and coherence, and emotional needs responsiveness. Four coaching modes---emotional support, practical education, policy navigation, and mixed---are governed by three pipeline invariants and classified through a hybrid LLM-and-heuristic intent router. The system is evaluated through automated pillar test matrices, latency and load benchmarks, and audit-trail validation, with Swiss informal caregivers as the pilot use case.

\begin{table}[H]
\centering\small
\caption{Research questions}
\begin{tabularx}{\textwidth}{@{}lX@{}}
\toprule
RQ & Focus \\
\midrule
RQ1 & Continuous OCEAN inference with EMA: valid, temporally stable trait estimates? \\
RQ2 & Personality modulation improves emotional tone, relevance, emotional-needs responsiveness vs.\ non-adaptive baselines? \\
RQ3 & RAG-enhanced policy navigation: full citation coverage, no critical grounding errors? \\
RQ4 & Hybrid LLM router correctly classifies intent across four coaching modes with acceptable latency? \\
RQ5 & System achieves acceptable reliability under load with full audit coverage? \\
\bottomrule
\end{tabularx}
\end{table}

Chapter~\ref{ch:related} surveys related work. Chapter~\ref{ch:system} describes the system design addressing all five RQs. Chapter~\ref{ch:eval} specifies the evaluation methodology. Chapter~\ref{ch:results} reports results (RQ1 completed; RQ2--5 pending). Chapter~\ref{ch:discussion} interprets findings and discusses limitations.

\chapter{Related Work}
\label{ch:related}

\section{Digital Coaching Systems and Personalisation Mechanisms}

Digital coaching has emerged as a scalable intervention format, with the AI-driven mental health app market reaching \$1.45 billion in 2024 \citep{GrandView2024}. Three architectural models dominate current deployments, each with characteristic personalisation capabilities and limitations.

\textbf{Rule-based systems} such as Woebot \citep{Fitzpatrick2017} deliver CBT through pre-programmed decision trees, with mood-adaptive branching selecting exercise modules based on keyword-classified emotional state. Woebot's RCT demonstrated significant depression symptom reductions versus an information-only control \citep{Fitzpatrick2017}, validating the therapeutic value of structured digital coaching. However, two users with identical mood states but different personality profiles receive identical content and delivery style.

\textbf{Hybrid NLP systems} such as Wysa \citep{Inkster2018} combine rule-based therapeutic modules with free-text NLP routing, crisis keyword detection, and working-alliance development through conversational rapport-building \citep{Beatty2022}. Personalisation is richer but still operates at the level of transient state and conversational context---the system adapts to what the user says and feels now, not to who they are across sessions.

\textbf{LLM-based memory systems} store user statements as vector embeddings for cross-session retrieval \citep{Li2024}. Comprehensive surveys identify three technical levels: prompting with personalised context, fine-tuning with personalised adapters, and alignment with personalised preferences \citep{Liu2025, Wang2026}.\footnote{\citet{Liu2025} and \citet{Wang2026} are preprints; cited as of manuscript submission (March 2026).} These systems enable impressive conversational continuity but personalise \emph{content recall} rather than \emph{coaching style}. A memory system knows what the user said; it does not know who the user is.

\textbf{The personality gap.} Across all three models, no deployed coaching system models the stable individual differences that predict how users respond to different coaching styles. Big Five traits moderate coaching outcomes---Openness predicts benefit from exploratory approaches, Conscientiousness from structured goal-setting, Neuroticism from emotional scaffolding \citep{deHaan2013, Jones2016}; individual-level factors similarly moderate performance improvement in structured feedback interventions \citep{Smither2005}---yet no digital system has operationalised this evidence. The cost of this gap is not merely suboptimal outcomes but measurably higher attrition: personality-blind systems produce emotionally incongruent responses that high-Neuroticism users experience as invalidating, amplifying rather than buffering the negative affect that originally drove them to seek support \citep{Eysenbach2005}. Digital health attrition research confirms this mechanism: interventions incorporating emotional validation and personalised encouragement sustain engagement for users in high-stress states, whilst those delivering generic content to emotionally activated users accelerate dropout \citep{Beatty2022}.

Critically, modern personality science treats traits not as fixed labels applied once but as \emph{distributions of states} that activate differentially across situational contexts: extraversion increases in social settings, neuroticism spikes under stress, and conscientiousness rises during task-oriented contexts \citep{Fleeson2001}. Experience sampling studies confirm that momentary personality states predict user outcomes---including engagement, well-being, and intervention response---independently of static trait scores \citep{Fleeson2009}, and that effective personalisation therefore requires per-turn dynamic inference rather than a single baseline assessment. No existing digital coaching system performs continuous within-session personality inference.

\textbf{Personality-aware dialogue.} The technical foundations for bridging this gap exist. \citet{Yarkoni2010} demonstrated systematic associations between word-use patterns and personality dimensions; \citet{Park2015} achieved personality prediction from social media text. In dialogue generation, \citet{Mairesse2010} showed that personality-conditioned language improves perceived naturalness; \citet{Hu2018} extended this to joint modelling of linguistic style and personality adaptation in NLG. Subsequent work explored persona-based dialogue models, including the PersonaChat dataset and framework \citep{Zhang2018, Li2016, Zheng2019}. \citet{Devdas2025} advanced continuous Big Five detection with motivational regulation, reporting improvements in emotional support quality. The Zurich Model of Social Motivation \citep{Quirin2022, Quirin2023} provides a theoretical bridge from OCEAN to coaching behaviour through three motivational systems---security (Neuroticism), arousal (Openness, Extraversion), and affiliation (Agreeableness)---with Conscientiousness as a structural modifier.

Two deployment problems remain unsolved. First, \emph{temporal stability}: persona-aware generation conditions on static profiles \citep{Li2016, Majumder2020} without modelling how personality estimates should evolve and stabilise across turns. The field confirms that ``traditional approaches often rely on static personality profiles,'' with dynamic representation identified as an open problem \citep{Zheng2019}. Recent systems address \emph{within-session} personality dynamism---COSPLAY \citep{Xu2022} represents self- and partner personas as evolving concept sets; the Dual Enhanced Network \citep{Zhu2024} jointly encodes long-term traits and short-term states for detection; \citet{Wen2024b} formalises personality as a transition weight in per-turn emotion generation---but none address \emph{cross-session stabilisation} of noisy per-turn detections into persistent persona profiles. A further open problem is identity drift: LLMs conditioned on personality profiles exhibit deteriorating Rank-Order stability under prolonged interaction, with core traits increasingly contradicted as conversations lengthen \citep{Kovac2024}. All dynamic approaches require model training on annotated dialogue corpora unavailable for Swiss caregiver scenarios. Second, \emph{classification to behaviour}: adaptive dialogue management has approached this through reinforcement learning, meta-learning, and multi-task learning \citep{denHengst2019, Lee2021}. RL-based personalisation---including persona-consistent offline RL \citep{Shea2023}---requires per-domain reward function design and thousands of training interactions before an initial policy is ready. Meta-learning approaches---such as personalising dialogue agents via MAML-style adaptation \citep{Lin2019}---achieve few-shot personalisation from dialogue history but require a meta-training phase across multiple users, making them unsuitable for cold-start deployment with populations (Swiss caregivers) for whom no prior interaction corpus exists. Both families of methods also produce policies difficult to audit at the level of individual decisions. No prior system has combined continuous trait detection, cross-session temporal stabilisation, and formal directive-based modulation in an operational coaching pipeline without training data. The Zurich Model operationalisation in this work---mapping OCEAN traits to behavioural directives through the security, arousal, and affiliation systems---is absent from the prior dialogue literature.

\section{Retrieval-Augmented Generation and Grounding}

Retrieval-augmented generation (RAG) addresses the well-documented tendency of LLMs to generate plausible but factually unsupported content---hallucination \citep{Ji2023}. \citet{Lewis2020} introduced the RAG framework, demonstrating that conditioning generation on retrieved documents improves factual accuracy across knowledge-intensive tasks. Hybrid retrieval strategies---combining dense vector similarity and sparse lexical matching---achieve superior recall compared to either method alone \citep{Karpukhin2020, Ma2023}.

In health-adjacent and policy-guidance domains, grounding is a safety requirement rather than a quality improvement. Incorrect guidance about benefit eligibility or application procedures can cause material harm to vulnerable populations. Multi-stage grounding verification---extracting claims, checking evidence overlap, and flagging or replacing unsupported assertions---has emerged as a necessary complement to retrieval \citep{Thoppilan2022, Gao2023, Singhal2023}. A comprehensive survey of hallucination mitigation techniques identifies post-generation verification as the most reliable approach for deployment contexts where retrieval quality cannot be guaranteed \citep{Tonmoy2024}. Within post-generation verification, Chain-of-Verification (CoVe) establishes the principle that decomposing a response into independently verifiable claims---rather than assessing the whole response at once---prevents undetected hallucination through evidence masking, reducing unverified claim rates by up to 28\% on open-domain factual tasks \citep{Dhuliawala2023}. Fine-grained claim-level detection further demonstrates that atomic sentence decomposition substantially outperforms response-level consistency scoring \citep{Mishra2024}.

Knowledge graph (KG) integration has been explored for enriching persona consistency---encoding persona attributes, interests, and belief structures to ensure character-coherent responses \citep{Majumder2020}. KG-based persona consistency is, however, distinct from factual policy grounding: the former ensures the system's persona remains internally coherent, whilst the latter enforces that specific policy claims trace to authoritative sources. No prior system has defined formal, enforced invariants that simultaneously require per-turn personality detection, conditional evidence retrieval, and citation-mandatory grounding. This combination is an architectural contribution distinct from either KG persona consistency or standard RAG deployment.

Big5Loop integrates RAG with personality-aware coaching through explicit architectural separation: retrieval and grounding verification govern \emph{what} the system says, whilst personality modulation governs \emph{how} it says it. Three enforced pipeline invariants maintain this separation. To our knowledge, this is the first system to treat personality modulation and factual grounding as simultaneously required but architecturally separated pipeline properties.

\section{Digital Support for Informal Caregivers}

Informal caregivers experience chronic stress, emotional exhaustion, and social isolation at rates significantly exceeding the general population \citep{Schulz2008, Pinquart2003}. Their stressor profile spans four intersecting domains: physical demands from direct care tasks and sleep disruption; emotional burden from witnessing decline and managing behavioural symptoms such as verbal agitation (present in 40--80\% of dementia care recipients); financial strain from reduced employment and increased care costs; and social isolation from constrained participation and strained family relationships \citep{Brodaty2009}. Digital interventions have shown promise---\citet{Mooney2024} reported burden reductions of up to 38\% at 8 weeks in a systematic review and meta-analysis---but are severely constrained by a personalisation gap that no existing tool has addressed.

Systematic reviews of digital caregiver interventions reveal a consistent pattern: high attrition rates, with many caregivers disengaging after only a few interactions, and limited efficacy in studies that do not account for individual differences. Research on digital health attrition identifies six primary driver categories---usability, perceived value, technical barriers, lack of personalisation, emotional and psychological factors, and contextual factors such as disease burden and demographics. For caregiver populations, the emotional/psychological and personalisation categories are dominant: caregivers contact digital tools precisely during acute stress episodes, and a generic system that delivers the same content regardless of the caregiver's current state provides a response that is emotionally incongruent---experienced not merely as irrelevant but as invalidating, which amplifies negative affect and actively accelerates disengagement rather than supporting it. Conversely, digital interventions incorporating emotional validation, positive reinforcement, and personalised encouragement sustain engagement for users in high-stress states.

Personality traits are significant moderators of both engagement and outcome. Neuroticism predicts higher attrition: high-Neuroticism individuals initially seek digital tools under acute stress but disengage rapidly when tools fail to provide emotionally congruent responses---the causal mechanism is that uniform delivery is interpreted as dismissal, amplifying rather than buffering emotional arousal. Conscientiousness predicts sustained engagement and adherence to structured features. Agreeableness predicts responsiveness to empathetic conversational agents. A single generic intervention design systematically fails the users who need it most: high-Neuroticism caregivers would benefit from emotionally calibrated support but abandon systems that deliver uniform, non-responsive responses. No existing caregiver tool integrates personality assessment with adaptive response style.

The most directly comparable prior system is the Caregiving Language Model (CaLM) \citep{Parmanto2024}, which uses a RAG framework combined with fine-tuned small language models (LLaMA-2 and Falcon, 7B parameters) grounded on a caregiving knowledge base for family caregivers of individuals with Alzheimer's Disease and Related Dementias. CaLM's central finding---that fine-tuned small models with RAG outperform GPT-3.5 on reference-grounded caregiving responses---establishes that factual grounding is both feasible and necessary in the caregiving domain. Big5Loop accepts this premise and treats grounding as a solved baseline problem. The open problem CaLM does not address is the \emph{how}: CaLM delivers the same response style to every caregiver regardless of trait profile, providing no mechanism for adapting emotional tone, structure, or pacing to the individual receiving guidance. The practical cost of this limitation is that a high-Neuroticism caregiver and a high-Conscientiousness caregiver receive identically framed responses despite requiring structurally different communication. Big5Loop's architectural contribution is therefore not an alternative grounding strategy but an additional layer that CaLM makes no attempt to provide: personality-modulated delivery enforced through formal pipeline invariants that simultaneously guarantee factual integrity and trait-adapted style.

The Swiss context adds a domain-specific dimension that makes this combination particularly necessary. Switzerland's federalised healthcare system produces canton-specific variations in policy eligibility and service availability. Caregivers must navigate IV/Invalidenversicherung, AHV, Ergänzungsleistungen, cantonal Spitex services, and Hilflosenentschädigung---a fragmented landscape demanding both emotional resilience to engage with complexity and accurate factual guidance to act on it. Social robot interventions such as ElliQ \citep{Broadbent2024} and evaluation frameworks such as FEEL \citep{Zhang2024} improve scalability but lack personality adaptation and Swiss-specific knowledge grounding. CBT-based coaching apps address emotional wellbeing through therapeutic exercises but do not deliver grounded policy guidance.

This architectural requirement---personality-modulated, citation-grounded, and fully anonymous---aligns with empirically documented barriers to formal help-seeking among Swiss informal caregivers. A qualitative study of family caregivers of people with dementia in Ticino (southern Switzerland) identified four primary barriers to seeking help: a high level of burden that left caregivers too overwhelmed to navigate services; a strong sense of duty framing care as a personal and family responsibility incompatible with external help-seeking; fear of being misunderstood or having difficulties minimised by those without direct caregiving experience; and difficulty accessing and navigating fragmented local services \citep{Messina2022}. The same study documented the paradox these barriers generate: caregivers who reported high burden and clear need for support simultaneously reported reluctance to actively seek it. Discrete, anonymous digital coaching directly addresses two of these barriers---it requires no disclosure of family difficulties to an institution, and it does not depend on caregivers navigating the fragmented service landscape that participants described as disorienting and taxing. The combination makes conversational AI a particularly well-matched support modality for this population.

No existing caregiver support tool integrates transient state tracking, conversation memory, stable personality modelling, and citation-grounded domain knowledge in a single adaptive system. Comprehensive surveys of the personality-aware digital coaching field---covering hundreds of papers on personality representation, adaptive dialogue, and evaluation methodology---report no caregiver-facing implementations, no caregiver-specific datasets, and no systems designed for Swiss policy contexts. The present study addresses a domain that is both practically urgent and academically unexplored.

\chapter{System Design}
\label{ch:system}

\section{Architecture and Pipeline Invariants}

Big5Loop is implemented as a three-tier architecture: a Next.js API route handles client requests and orchestrates pre- and post-processing; an N8N v2 workflow engine (version 2.12.0) executes the coaching pipeline; and PostgreSQL 16 with pgvector provides persistent storage for sessions, personality states, policy evidence, and audit records. Meta Llama 3 70B Instruct \citep{Dubey2024} (\texttt{meta/llama3-70b-instruct}, 70-billion-parameter instruction-tuned model), accessed via the NVIDIA API, serves as the language model for personality detection, response generation, and intent routing.

Each conversational turn proceeds through three phases. \emph{API-side pre-processing} performs hybrid LLM intent routing, resolves session isolation and route keys, and conditionally executes vector and full-text retrieval to prefetch policy evidence. The \emph{N8N workflow} then executes a partially parallelised pipeline: input normalisation and context extraction; loading of prior personality states from the database; then---in parallel---OCEAN trait detection with EMA smoothing \emph{and} policy retrieval via full-text search against a GIN-indexed \texttt{policy\_chunks} table; regulation through trait-to-directive mapping; merging of detection and retrieval results; directive-conditioned response generation; grounding verification through claim extraction and evidence matching; quality verification and safety checking; and response formatting with persistence of session, turn, personality, and metrics data. The parallelisation of detection and retrieval reduces end-to-end latency by eliminating the sequential dependency between these two stages (see \S\ref{sec:perf_opt}). A detection cache (in-memory, keyed on the first 300 characters of the user message, 5-minute TTL) further reduces latency on repeated messages by bypassing the LLM API entirely. \emph{API-side post-processing} applies mode override and confidence propagation, performs optional LLM regeneration for heuristic-generated content, overlays citations, executes the grounding invariant check, and writes dual audit records.

Three pipeline invariants govern every turn regardless of coaching mode. These are enforced architectural rules, not configurable options: making them invariant rather than configurable prevents silent degradation under operational conditions where retrieval latency or API availability might otherwise tempt downstream shortcuts, and ensures that no deployment configuration can accidentally disable safety-critical behaviour.

\begin{table}[H]
\centering\small
\caption{Pipeline invariants}
\begin{tabularx}{\textwidth}{@{}lXX@{}}
\toprule
Invariant & Rule & Enforcement \\
\midrule
Universal personality & OCEAN detection and regulation run for every turn in every mode & N8N executes Detection $\rightarrow$ EMA $\rightarrow$ Regulation before mode-specific branching; API flags \texttt{ocean\_detection\_skipped} or \texttt{regulation\_skipped} on violation \\
Conditional retrieval & Evidence retrieval runs only when factual content is needed & API-side gating function evaluates resolved mode, keyword signals, and policy opt-out phrases before initiating retrieval \\
Mandatory grounding & Policy and mixed responses must include citations & N8N grounding verifier checks claim-evidence overlap; API checks citation count post-response; zero citations triggers a grounding warning \\
\bottomrule
\end{tabularx}
\end{table}

The first invariant ensures every response is personality-informed regardless of coaching mode. The second prevents unnecessary retrieval latency on emotional or educational turns. The third guarantees no policy claim reaches the user without evidence support.

The four coaching modes define how the system balances personality modulation and factual grounding:

\begin{table}[H]
\centering\small
\caption{Coaching modes}
\begin{tabularx}{\textwidth}{@{}lXXX@{}}
\toprule
Mode & Personality modulation & Retrieval & Grounding \\
\midrule
Emotional support & Full (tone, warmth, pacing) & Skipped & Not required \\
Practical education & Full (structure, detail, style) & Conditional & If citations present \\
Policy navigation & Presentation style only & Always & Mandatory \\
Mixed & Full modulation + presentation style & Always & Mandatory for policy segment \\
\bottomrule
\end{tabularx}
\end{table}

The mixed mode is architecturally significant: when a caregiver expresses emotional distress and asks a policy question in the same message, the system composes a response with both an emotionally supportive segment and a citation-grounded policy segment, applying personality modulation throughout whilst maintaining factual integrity in the policy portion.

\begin{figure}[H]
\centering
\includegraphics[width=0.95\textwidth]{figure1_system_architecture}
\caption{Big5Loop three-tier architecture: Next.js API route, N8N v2 orchestration pipeline with parallel detection and retrieval branches (detect/retrieve $\rightarrow$ merge $\rightarrow$ regulate $\rightarrow$ generate $\rightarrow$ verify $\rightarrow$ persist $\rightarrow$ audit), and PostgreSQL/pgvector storage layer. Personality detection and EMA stabilisation run on every turn (invariant 1); retrieval activates conditionally (invariant 2); all policy content is citation-grounded (invariant 3). A detection cache bypasses the LLM API for repeated messages.}
\end{figure}

\section{Personality Detection and Stabilisation}

Personality detection, temporal stabilisation, and trait-to-directive modulation form a tightly coupled subsystem that runs on every turn, enforcing the universal personality invariant.

\textbf{Detection.} Each turn, the system infers current Big Five trait expression using Meta Llama 3 70B Instruct \citep{Dubey2024} at temperature 0.1 via few-shot prompting. The prompt comprises three components: (1) a calibrated system message defining trait-specific linguistic markers for each OCEAN dimension; (2) seven domain-representative few-shot examples spanning the full OCEAN range, each pairing a user utterance with the expected JSON trait scores; and (3) the user's current message. The model returns per-trait values in $[-1.0, +1.0]$ with per-trait confidence set to 0.8 for API-sourced detections. The few-shot calibration is critical: zero-shot prompting produced trait inversions (e.g., negative Conscientiousness correlation) because the model conflated topic sentiment with speaker personality style. A heuristic keyword fallback activates when the LLM is unavailable, providing approximate trait estimates at confidence 0.45.

\textbf{EMA stabilisation.} Raw per-turn estimates are noisy. The system addresses this through a hierarchical state estimation pipeline: LLM observations $\rightarrow$ session state estimation $\rightarrow$ long-term trait estimation $\rightarrow$ policy activation. This four-layer architecture operates across two time scales, mirroring experience sampling methodology (ESM/EMA) in personality science: within-person \emph{personality states} fluctuate moment-to-moment, whilst between-person \emph{personality traits} represent stable long-term baselines \citep{Fleeson2001, Nesselroade1991}.

\emph{Layer 1 --- Observation.} Each turn produces independent per-trait noisy observations with confidence scores: $(\hat{\mathbf{T}}^{(n)}, \mathbf{c}^{(n)}) = f_{\text{LLM}}(u^{(n)}, H^{(n-6:n-1)})$, $\hat{\mathbf{T}}^{(n)} \in [-1,1]^5$, where $u^{(n)}$ is the user message at turn $n$, $H$ is the 6-turn dialogue history, and each trait dimension $i \in \{\text{O,C,E,A,N}\}$ has confidence $c_i^{(n)} \in [0,1]$.

\emph{Layer 2 --- In-session state estimator.} Per-trait confidence gating prevents low-quality observations from corrupting the session estimate. For each trait dimension $i$ independently:
\begin{align}
T_{\text{sess},i}^{(n)} &= \begin{cases}
\alpha_s \cdot \hat{T}_i^{(n)} + (1-\alpha_s)\cdot T_{\text{sess},i}^{(n-1)} & \text{if } c_i^{(n)} \geq \tau_c \\
T_{\text{sess},i}^{(n-1)} & \text{otherwise}
\end{cases}
\end{align}
with $\alpha_s = 0.3$ and confidence gate $\tau_c = 0.4$. The value $\alpha_s = 0.3$ was selected to balance responsiveness to new evidence with resistance to single-turn noise; sensitivity analysis across $\alpha_s \in \{0.1, 0.2, 0.3, 0.4, 0.5\}$ shows convergence to within $\pm 0.05$ of the underlying trait signal within 6--8 turns. The threshold $\tau_c = 0.4$ was set to reject heuristic-fallback detections (confidence 0.45) while accepting LLM-sourced detections (confidence 0.8), ensuring that only higher-quality observations update the session estimate. The session is initialised from the prior persona profile: $T_{\text{sess},i}^{(0)} = P_i^{(k-1)}$.

\emph{Layer 3 --- Cross-session trait estimator.} After each session $k$ concludes: $P_i^{(k)} = \alpha_p \cdot T_{\text{sess},i}^{(\text{final})} + (1-\alpha_p)\cdot P_i^{(k-1)}$ with $\alpha_p = 0.2$. This value reflects the assumption that stable traits require approximately 5 sessions to update meaningfully, giving an effective memory horizon of $1/\alpha_p = 5$ sessions.

This two-level EMA structure is a fixed-gain approximation to a Kalman filter: under the assumption that the latent trait evolves as a random walk with constant process noise, the Kalman gain converges toward a constant value and the recursive update reduces to exponential moving average smoothing. The confidence gate $\tau_c$ is a computationally efficient hard approximation to continuous noise weighting.

\emph{Layer 4 --- Directive activation.} Stabilised session estimates trigger behavioural directives when trait magnitude exceeds $\theta_T = 0.15$: $\mathcal{D}^{(n)} = \{ d_i : |T_{\text{sess},i}^{(n)}| > \theta_T \}$.

\begin{table}[H]
\centering\small
\caption{Trait-to-directive mapping (summary)}
\begin{tabularx}{\textwidth}{@{}lXc@{}}
\toprule
Trait direction & Behavioural directive & Activation \\
\midrule
High Neuroticism & Offer extra comfort; acknowledge anxieties & $T_{\text{sess},N} > 0.15$ \\
Low Neuroticism & Reassure stability and confidence & $T_{\text{sess},N} < -0.15$ \\
High Agreeableness & Show warmth, empathy, collaboration & $T_{\text{sess},A} > 0.15$ \\
High Conscientiousness & Provide organised, structured guidance & $T_{\text{sess},C} > 0.15$ \\
High Openness & Invite exploration and novelty & $T_{\text{sess},O} > 0.15$ \\
High Extraversion & Energetic, sociable tone & $T_{\text{sess},E} > 0.15$ \\
\bottomrule
\end{tabularx}
\end{table}

When no trait exceeds the activation threshold, default directives apply: warm, attentive, and supportive with open-ended questions. The table above shows high-direction variants only; the full 10-row mapping including low-direction variants for all five traits is in Appendix~\ref{app:directives}. The directive mapping operationalises documented empirical patterns (e.g., high-Neuroticism caregivers interpret demands as threats; high-Conscientiousness employ structured problem-solving) \citep{Carver1989, Lazarus1984, Penley2002}. Big5Loop implements this through explicit trait-to-directive rules rather than learned discriminators, preserving interpretability and auditability.

\begin{figure}[H]
\centering
\includegraphics[width=0.85\textwidth]{figure2_ema_convergence}
\caption{EMA trait stabilisation: raw per-turn Neuroticism detections (volatile) versus EMA-smoothed in-session estimate and cross-session trait baseline across 8 turns. Confidence gating at turns with $c_i^{(n)} < \tau_c = 0.4$ preserves the prior estimate. The EMA estimate converges within 6 turns to a stable neighbourhood.}
\end{figure}

\section{Hybrid Intent Router}

Classifying the user's coaching need determines which pipeline branches activate---particularly whether evidence retrieval and mandatory grounding are triggered. Context-aware NLU systems that incorporate dialogue history consistently outperform single-turn classifiers \citep{Gupta2019b}; however, large pretrained models are prone to overconfidence on out-of-domain inputs, motivating confidence gating as a rejection mechanism \citep{Wang2024}. Big5Loop's router spans both task-oriented intents (policy navigation, practical education) and open-domain social intents (emotional support, mixed)---a combination absent from prior intent classification systems \citep{Liu2024}.

The system uses a five-layer hybrid router: (1) \textbf{Explicit requested mode}---frontend-specified modes used directly with confidence 1.0; (2) \textbf{LLM router}---Llama 3 70B Instruct at temperature 0 returns structured JSON (mode, confidence, retrieval flag, routing reason); results cached per message-context key for 5 minutes; conversation-history context (prior 3 turns) included; (3) \textbf{Confidence gate}---confidence below 0.5 discards the LLM decision and falls through to heuristics; (4) \textbf{Keyword heuristic fallback}---curated term lists (26 policy, 17 education, 15 emotional terms) with conversation-history inheritance for short follow-up messages; (5) \textbf{Hard safety override}---policy-related keywords force retrieval activation regardless of routing source. The N8N V2 workflow consumes the API routing decision directly. Every response carries routing metadata (classification source, routing reason, confidence score, LLM latency).

\section{Retrieval and Grounding Verification}

Once the router resolves the coaching mode, the retrieval and grounding subsystem activates conditionally: emotional support turns bypass it entirely, while policy and mixed turns trigger the full retrieval and verification pipeline described below.

\textbf{Retrieval.} When the resolved mode requires evidence, the system executes hybrid retrieval. The primary path uses the NVIDIA embedding API (model \texttt{nv-embedqa-e5-v5}, 1024-dimensional vectors) for cosine similarity against pre-embedded policy chunks in PostgreSQL/pgvector, returning the top 5 results above a 0.30 similarity threshold. When vector search returns fewer than 2 results, PostgreSQL full-text search (\texttt{websearch\_to\_tsquery}) supplements and deduplicates. The policy corpus covers selected Swiss social insurance domains (IV eligibility, Zurich supplementary benefits, care subsidies, Hilflosenentschädigung). Personality modulation does not alter retrieved content---it governs only how that content is presented.

\textbf{Grounding verification.} A three-stage pipeline ensures no policy claim reaches the user without evidence support. \emph{Stage 1:} The N8N Grounding Verifier node extracts policy claims via regex patterns and checks keyword overlap against retrieved evidence. A grounding ratio below 0.5 triggers full safe-fallback replacement; 0.5--0.8 produces a degraded response replacing ungrounded claims with ``[information pending verification]''; $\geq 0.8$ passes. \emph{Stage 2:} The Verification and Refinement node scores overall adherence (session validity, question count, word count, banned claim phrases); scores below 0.5 trigger a full block. \emph{Stage 3:} The API performs a post-response citation count check for policy and mixed mode responses; zero citations generates a \texttt{warn\_no\_citations} audit warning---a final backstop that no prior RAG system implements as a formal pipeline invariant.

\section{Performance Optimisation}
\label{sec:perf_opt}

End-to-end latency is dominated by external LLM API calls (detection and generation), which together account for over 95\% of turn processing time. Three optimisations were implemented and evaluated via A/B comparison between a sequential baseline and the optimised parallel pipeline.

\textbf{Parallel detection and retrieval.} In the original sequential pipeline, personality detection completed before policy retrieval began. Since the detection result (OCEAN scores) is not required for the initial full-text retrieval query---which uses only the user's message text---these two stages can execute concurrently. The N8N workflow was restructured to branch from the ``Merge Previous State'' node into two parallel paths: one executing ``Zurich Model Detection (EMA)'' $\rightarrow$ ``Enhanced Regulation,'' and the other executing ``Policy Retrieval (PostgreSQL).'' A merge node recombines results before generation. PostgreSQL full-text search uses a GIN index on \texttt{to\_tsvector('english', content || ' ' || title)} over the \texttt{policy\_chunks} table, ensuring sub-millisecond retrieval latency.

\textbf{Detection cache.} An in-memory cache (keyed on the first 300 characters of the user message, 5-minute TTL, stored via N8N's \texttt{\$getWorkflowStaticData('global')}) stores previously computed OCEAN scores. Cache hits bypass the NVIDIA API entirely, reducing detection latency from ${\sim}4$--5\,s to $<50$\,ms. The cache is particularly effective for repeated or near-identical messages within a session.

\textbf{Reduced API parameters.} Detection timeout was reduced from 30\,s to 12\,s and \texttt{max\_tokens} from 100 to 50 (sufficient for the compact JSON output); generation timeout was reduced from 45\,s to 20\,s and maximum tokens were tightened per model tier (light: 80, medium: 120, heavy: 150).

\begin{table}[H]
\centering\small
\caption{Pipeline latency comparison: sequential vs.\ parallel architecture (5-request average)}
\label{tab:perf}
\begin{tabular}{@{}lccc@{}}
\toprule
Metric & Sequential & Parallel (fresh) & Parallel (cached) \\
\midrule
Detection (ms) & $4500\pm800$ & $4200\pm700$ & $<50$ \\
Retrieval (ms) & $50\pm20$ & \multicolumn{2}{c}{runs in parallel with detection} \\
Generation (ms) & $3200\pm600$ & $2800\pm500$ & $2800\pm500$ \\
End-to-end (ms) & $9500\pm1200$ & $8100\pm900$ & $3500\pm600$ \\
\midrule
Improvement vs.\ sequential & --- & ${\sim}15\%$ & ${\sim}63\%$ \\
\bottomrule
\end{tabular}
\end{table}

The parallel architecture yields a ${\sim}15\%$ improvement for fresh requests (where detection must call the LLM API) and a ${\sim}63\%$ improvement for cached requests (where detection is near-instantaneous). The bottleneck remains the external LLM API; further latency reduction requires either model distillation or local inference.

\section{Infrastructure, Audit, and Reproducibility}

Big5Loop is deployed via Docker Compose: a Next.js 14 frontend (Turbopack-accelerated dev server) with Zustand state management; N8N v2.12.0 self-hosted with JavaScript task runners; PostgreSQL 16 with pgvector across 11 tables (sessions, turns, personality states, policy evidence, metrics, audit logs, users, feedback, memory embeddings, user personality profiles, policy chunks); and Meta Llama 3 70B Instruct via NVIDIA API (\texttt{meta/llama3-70b-instruct}). The N8N upgrade from v1 to v2 required migrating in-node state persistence from \texttt{globalThis} to \texttt{\$getWorkflowStaticData('global')} and replacing the \texttt{update:workflow --active} CLI with \texttt{publish:workflow}. Inter-module contracts are validated using Zod schemas in a shared contracts package (\texttt{@big5loop/contracts}); user authentication uses HS256 JWT tokens in HTTP-only cookies.

Every turn produces a dual audit record---JSONL log and PostgreSQL \texttt{audit\_log} row---containing: request ID, session ID, turn index, coaching mode, pipeline status, routing metadata, personality state (OCEAN scores only, no PII), retrieval IDs, citation count, verifier status, SHA-256 input hash, and turn latency. PII redaction precedes all logging.

Reproducibility is supported by: fixed model versions and parameters (Llama 3 70B Instruct, \texttt{meta/llama3-70b-instruct}, EMA $\alpha_s = 0.3$); versioned N8N workflow files (V-numbering convention); deterministic JSON contracts between all modules; containerised single-command deployment; and benchmark fixtures stored as versioned JSON enabling automated regression testing.

\chapter{Evaluation}
\label{ch:eval}

\section{Study Design and Validation Streams}

This study follows a pilot system design and evaluation approach. Four validation streams provide converging evidence:

\begin{table}[H]
\centering\small
\caption{Validation streams}
\begin{tabularx}{\textwidth}{@{}lXX@{}}
\toprule
Stream & Sample & What it validates \\
\midrule
PERSONAGE benchmark & 320 human-annotated utterances with ground-truth OCEAN \citep{Mairesse2011} & External trait detection validity (Pearson $r$, MAE) \\
Automated pillar test matrices & 130+ cases across 4 modes & Mode classification, citation presence, grounding safety, latency gates \\
Benchmark and load tests & 10 expansion + 4 latency + 100 concurrent sessions & Performance gates, reliability under load \\
Expert pilot sessions & $n = 5$--8 domain experts (Spitex coordinators, geriatricians, home care nurses) & Usability, appropriateness, ecological validity \\
Simulated evaluation & $N \geq 250$ synthetic conversations across personality profiles and stress scenarios & Coaching quality compared against baselines \\
\bottomrule
\end{tabularx}
\end{table}

The PERSONAGE benchmark \citep{Mairesse2011} provides external validation of trait detection accuracy against human-annotated ground truth. PERSONAGE contains restaurant-recommendation utterances systematically generated to express controlled Big Five personality profiles, with each utterance rated by human judges on all five OCEAN dimensions (1--7 Likert scale, normalised to $[-1, +1]$). The corpus contains 320 utterances before final filtering: 160 from predefined personality conditions and 160 from randomised parameter combinations. After filtering for complete ground-truth and detected OCEAN vectors, 318 valid samples remain (reported in \S\ref{ch:results}). Real caregiver validation ($n = 20$--30) is fully specified as a future protocol, deferred due to recruitment feasibility constraints.

\section{Baselines}

Four conditions isolate the contribution of each personalisation layer:

\begin{table}[H]
\centering\small
\caption{Experimental conditions and personalisation layers}
\begin{tabularx}{\textwidth}{@{}lXX@{}}
\toprule
Condition & Personalisation layers active & What it isolates \\
\midrule
Personality-aware (full system) & Mood + memory + personality + RAG & Full personalised coaching \\
Generic non-adaptive & None (static responses) & Whether any personalisation matters \\
Memory-only & Memory + mood (no trait modelling) & Whether conversational continuity substitutes for personality modulation \\
Policy-only & RAG (no personality or memory) & Whether factual accuracy alone suffices \\
\bottomrule
\end{tabularx}
\end{table}

The memory-only versus full-system comparison is particularly informative: memory represents the current state-of-the-art in LLM personalisation, and any difference isolates the specific value of personality-trait modelling.

\section{Measures}

\begin{table}[H]
\centering\footnotesize
\caption{Evaluation metrics, targets, and RQ mapping}
\begin{tabularx}{\textwidth}{@{}lXXc@{}}
\toprule
Construct & Metric & Target & RQ \\
\midrule
Trait validity (external) & Pearson $r$ with PERSONAGE ground truth & $\geq 0.40$ (moderate); $\geq 0.75$ (strong) & RQ1 \\
Trait validity (internal) & Pearson $r$ with synthetic profiles & $\geq 0.75$ & RQ1 \\
Trait stability & Variance within 6+ turns & $< 0.05$ & RQ1 \\
Confidence calibration & High-confidence predictions aligned & $\geq 85\%$ & RQ1 \\
Emotional tone appropriateness & LLM evaluator (CoT, temp 0) + human expert & $\geq 4.0/5.0$ & RQ2 \\
Relevance \& coherence & LLM evaluator (CoT, temp 0) + human expert & $\geq 4.0/5.0$ & RQ2 \\
Emotional needs responsiveness & Improvement over baseline & $\geq 20\%$ & RQ2 \\
Citation coverage & Policy claims with sources & 100\% & RQ3 \\
Grounding errors & Critical hallucinations & 0 & RQ3 \\
Routing accuracy & Correct mode classification (benchmark) & 10/10 & RQ4 \\
LLM router latency & Per-turn classification time & Monitored & RQ4 \\
Load reliability & 100 concurrent sessions, 5xx rate & 0\% & RQ5 \\
End-to-end latency (p95) & Emotional / practical / policy / mixed & $\leq 4$ / 5 / 8 / 9 s & RQ5 \\
Audit coverage & Turns with complete audit records & 100\% & RQ5 \\
\bottomrule
\end{tabularx}
\end{table}

Traditional automatic metrics (BLEU, ROUGE, perplexity) are not used as primary evaluation metrics; they measure lexical overlap rather than personality fit or coaching quality \citep{Yeh2021}. The simulated evaluation uses an LLM evaluator (Llama 3 70B Instruct, temperature 0) scoring each response on three dimensions: emotional tone appropriateness, relevance and coherence, and emotional needs responsiveness. Three design choices address LLM evaluator risks \citep{Guo2023}: (1) chain-of-thought prompting---each evaluator response must provide a step-by-step rationale before assigning a score; (2) temperature 0---deterministic outputs; (3) fixed evaluation rubrics per dimension. Human expert gating ($\kappa \geq 0.75$ inter-rater agreement on a validation subset) serves as a final calibration layer.

\section{Procedure}

(1) Build and deploy the full coaching pipeline (Next.js + N8N + PostgreSQL + NVIDIA API). (2) Implement the hybrid LLM router with heuristic fallback and confidence gating. (3) Implement and verify the three pipeline invariants. (4) Execute automated pillar test matrices across all four coaching modes (130+ cases). (5) Execute benchmark expansion pack (10 cases) and latency benchmarks (4 cases). (6) Execute load test at 100 concurrent sessions. (7) Conduct expert pilot sessions; collect system logs and appropriateness ratings. (8) Generate $N \geq 250$ synthetic conversations across personality profiles and stress scenarios. (9) Score synthetic conversations using the automated LLM evaluator; validate on human expert subset (target $\kappa \geq 0.75$). (10) Compare personality-aware system against all three baseline conditions.

\textbf{Analysis plan.} RQ1: Pearson correlation with PERSONAGE ground truth (external validity); Pearson $r$ with synthetic profiles (internal validity); variance across turns. RQ2: Adaptive vs.\ baseline on three coaching-quality dimensions; Cohen's $d \geq 0.3$ target. RQ3: Policy claim audit; citation coverage count; error categorisation. RQ4: Routing accuracy across 4 modes; confidence distribution; LLM vs.\ heuristic comparison. RQ5: Load test pass rate; latency percentiles; audit coverage; contract compliance rate.

\section{Ethics}

This study complies with the Swiss Federal Act on Data Protection (FADP). All data are pseudonymised at collection; PII (e.g., email addresses, phone numbers) is removed by automated redaction before logging. The dual audit trail enables post-hoc transparency whilst maintaining privacy. A crisis protocol escalates automatically to professional resources when detected stress level reaches 3 or above on a 0--4 scale. The system maintains a non-clinical boundary: it generates no medical diagnoses, legal advice, or clinical treatment recommendations. Grounding verification functions as a safety mechanism, preventing fabricated policy claims from reaching vulnerable users.

\chapter{Results}
\label{ch:results}

\section{Personality Detection and Stabilisation (RQ1)}

\subsection{External Validation: PERSONAGE Benchmark}

To establish convergent validity of the personality detector against independently annotated ground truth, we evaluated the few-shot Llama 3 70B Instruct detection pipeline on the complete PERSONAGE corpus: 320 utterances with human-annotated ground-truth OCEAN ratings \citep{Mairesse2011}. The corpus spans all personality conditions: 160 predefined-parameter utterances (20 each for agree\_high, agree\_low, consc\_high, consc\_low, ems\_high, ems\_low, open\_high, open\_low) and 160 randomised-parameter utterances. Each utterance was processed in isolation (single-turn, no conversation history); the detector returned per-trait scores in $[-1.0, +1.0]$. After filtering for complete ground-truth and detected OCEAN vectors, 318 valid samples remain.

\begin{table}[H]
\centering\small
\caption{PERSONAGE trait detection results ($n = 318$, notebook run): Pearson $r$, $p$-value, MAE, sign agreement, and mean signed error (positive = overestimation)}
\begin{tabular}{@{}lccccc@{}}
\toprule
Trait & Pearson $r$ & $p$-value & MAE & Sign agree & Mean signed err \\
\midrule
Openness (O)          & 0.409 & $3.0 \times 10^{-14}$ & 0.318 & 64.5\% & $-$0.032 \\
Conscientiousness (C) & 0.425 & $2.2 \times 10^{-15}$ & 0.391 & 72.3\% & $-$0.053 \\
Extraversion (E)      & 0.390 & $5.4 \times 10^{-13}$ & 0.383 & 69.2\% & $-$0.045 \\
Agreeableness (A)     & 0.601 & $1.3 \times 10^{-32}$ & 0.285 & 79.2\% & $-$0.025 \\
Neuroticism (N)       & 0.526 & $5.0 \times 10^{-24}$ & 0.337 & 68.9\% & $+$0.075 \\
\textbf{Mean}         & \textbf{0.470} & & \textbf{0.343} & \textbf{70.8\%} & \\
\bottomrule
\end{tabular}
\end{table}

All five traits show statistically significant positive correlations (all $p < 5 \times 10^{-13}$). Agreeableness ($r = 0.60$) and Neuroticism ($r = 0.53$) show the strongest correlations, consistent with the salience of their linguistic markers in PERSONAGE: politeness hedges and collaborative phrasing (``would suggest'', ``you would like'', ``buddy'') reliably index high Agreeableness, while stuttering repetitions, hedging phrases (``I might be wrong'', ``err''), and exclamatory volatility markers reliably index high Neuroticism. Openness, Conscientiousness, and Extraversion yield moderate correlations ($r = 0.41$, $0.43$, $0.39$), reflecting greater cue overlap---vocabulary richness, formal register, and enthusiastic tone are harder to disentangle from single utterances without multi-turn context.

The mean signed error column reveals a systematic pattern: the detector slightly underestimates trait magnitude for four of five traits (O, C, E, A; mean signed errors $-0.025$ to $-0.053$), while modestly overestimating Neuroticism ($+0.075$). This asymmetry is consistent with the prompt's stronger negative-N cues (hedging, frustration markers) relative to the positive-N anchors. The underestimation bias is small relative to MAE and does not affect directional coaching decisions.

\textbf{Context and interpretation.} The mean $r = 0.47$ falls below the a priori convergent validity target $r \geq 0.75$. However: PERSONAGE utterances are short restaurant recommendations (mean $\sim$40 words), providing far less personality-informative text than multi-turn caregiver dialogues; the EMA stabilisation layer accumulates evidence across turns; and correlations are comparable to human inter-rater reliability for personality judgements from brief text, which \citet{Mairesse2011} report in the range $r = 0.40$--$0.65$ for the same PERSONAGE corpus. All correlations are consistently positive---the detector never inverts a trait's polarity. Zero-shot prompting produced negative Conscientiousness ($r = -0.40$) and near-zero Agreeableness ($r = -0.02$) because the model conflated topic-level sentiment with speaker personality style; the calibrated few-shot examples break this topic--personality conflation.

\subsection{Style-Aware Ensemble (Ablation)}

On a balanced 160-sample PERSONAGE subset with three prompt variants (standard few-shot, trait-first, anti-conflation), the style-aware ensemble yielded mean Pearson $r = 0.596$ compared to $r = 0.470$ on the full pipeline with the standard prompt---an improvement of +0.13. Conscientiousness and Neuroticism improve most markedly ($\Delta r = +0.13$ and $+0.10$ respectively). Agreeableness remains the best-detected trait across both configurations ($r = 0.60 \to 0.64$). Table~\ref{tab:ensemble} reports per-trait results for the isolated ensemble evaluation.

\begin{table}[H]
\centering\small
\caption{Isolated ensemble evaluation on balanced PERSONAGE subset ($n = 160$): Pearson $r$, MAE, and sign agreement}
\label{tab:ensemble}
\begin{tabular}{@{}lccc@{}}
\toprule
Trait & Pearson $r$ & MAE & Sign agree \\
\midrule
Openness (O) & 0.425 & 0.329 & 69.4\% \\
Conscientiousness (C) & 0.557 & 0.306 & 74.4\% \\
Extraversion (E) & 0.500 & 0.428 & 75.0\% \\
Agreeableness (A) & 0.644 & 0.361 & 77.5\% \\
Neuroticism (N) & 0.627 & 0.362 & 70.6\% \\
\textbf{Mean} & \textbf{0.596} & \textbf{0.357} & \textbf{73.4\%} \\
\bottomrule
\end{tabular}
\end{table}

\subsection{Visualisation of Detection Results}

Figure~\ref{fig:personage_scatter} shows per-trait scatter plots of detected versus ground-truth OCEAN scores for the full PERSONAGE corpus ($n = 318$). Figure~\ref{fig:personage_metrics} shows Pearson $r$ and MAE per trait. The PERSONAGE evaluation tests single-turn detection; in Big5Loop's operational context, EMA smoothing accumulates evidence across turns, and sensitivity analysis (\S3.2) demonstrates convergence to within $\pm 0.05$ of the underlying trait signal within 6--8 turns.

\begin{figure}[H]
\centering
\includegraphics[width=\textwidth]{personage_detected_vs_ground_truth}
\caption{PERSONAGE detection: Big5Loop-detected OCEAN scores versus human-annotated ground-truth scores, per trait ($n = 318$). Dashed lines indicate perfect agreement ($y = x$).}
\label{fig:personage_scatter}
\end{figure}

\begin{figure}[H]
\centering
\includegraphics[width=0.9\textwidth]{personage_metrics}
\caption{Per-trait Pearson $r$ and MAE for the full PERSONAGE corpus ($n = 318$). Agreeableness achieves the strongest detection ($r = 0.60$, MAE = 0.285); Extraversion is the weakest ($r = 0.39$, MAE = 0.383).}
\label{fig:personage_metrics}
\end{figure}

\section{Coaching Quality (RQ2)}

\begin{tcolorbox}[colback=yellow!8!white, colframe=orange!70!black, title=\textbf{Placeholder --- Coaching Quality Evaluation (RQ2)}, fonttitle=\bfseries, left=4pt, right=4pt, top=4pt, bottom=4pt]
\textit{Pending LLM-as-judge evaluation run. This section will report:}

\textbf{Experimental conditions} (four-way comparison):
\begin{enumerate}[leftmargin=*, nosep]
  \item \textbf{Adaptive Big5Loop} --- full pipeline: personality detection + EMA + directives + memory + RAG
  \item \textbf{Memory-only baseline} --- cross-session memory, no trait modulation (closest current SOTA)
  \item \textbf{Generic baseline} --- same LLM and retrieval; no personality, no memory
  \item \textbf{Policy-only baseline} --- RAG grounding without coaching or personality layer
\end{enumerate}

\textbf{Outcome dimensions} (5-point Likert, LLM-as-judge: Llama 3 70B Instruct, CoT, temperature 0, fixed rubric; human expert gating on 20-sample calibration set):
\begin{itemize}[leftmargin=*, nosep]
  \item \textbf{Emotional tone appropriateness (ETA):} Does the response style match the caregiver's inferred personality? (e.g., validating for high-N; structured for high-C; warm/exploratory for high-A)
  \item \textbf{Relevance and coherence (R\&C):} Is the response topically relevant, internally consistent, and appropriately scoped to the coaching mode?
  \item \textbf{Responsiveness to emotional needs (REN):} Does the response address the expressed emotional state appropriately (acknowledgment, reframing, validation)?
\end{itemize}

\textbf{Analysis plan:} Paired Wilcoxon tests between conditions; Cohen's $d$ effect sizes; breakdown by coaching mode (emotional support, practical education, policy navigation, mixed); failure-mode analysis for low-scoring adaptive responses. Target: adaptive $\geq 4.0/5.0$ on ETA and REN; $\geq 20\%$ improvement over generic on REN.

\medskip
\textbf{[Figure placeholder]} Grouped bar chart (or radar chart) --- adaptive vs.\ baselines across ETA, R\&C, REN, per mode and overall. Error bars = 95\% CI.

\medskip
\textbf{[Table placeholder]} Coaching quality scores per condition $\times$ dimension: mean $\pm$ SD, Cohen's $d$ vs.\ generic baseline, significance indicators.
\end{tcolorbox}

\section{Grounding and Routing (RQ3 + RQ4)}

\begin{tcolorbox}[colback=yellow!8!white, colframe=orange!70!black, title=\textbf{Placeholder --- Grounding \& Routing Evaluation (RQ3 + RQ4)}, fonttitle=\bfseries, left=4pt, right=4pt, top=4pt, bottom=4pt]
\textit{Pending pillar test matrix execution. This section will report:}

\textbf{Grounding (RQ3)} --- from the policy-mode pillar test matrix:
\begin{itemize}[leftmargin=*, nosep]
  \item \textbf{Citation coverage:} proportion of policy responses with $\geq 1$ source citation, by claim type (benefit eligibility, deadlines, monetary amounts, form requirements). Target: 100\%.
  \item \textbf{Verifier outcome distribution:} proportion of cases passing at $\geq 0.8$ (full pass), $0.5$--$0.8$ (degraded response), $< 0.5$ (safe-fallback replacement); Stage 2 block rate; Stage 3 \texttt{warn\_no\_citations} trigger rate.
  \item \textbf{Critical hallucinations:} count of ungrounded policy claims reaching the user in any test case. Target: 0.
  \item \textbf{Threshold calibration:} sensitivity of the 0.5/0.8 grounding ratio thresholds on the test set.
\end{itemize}

\textbf{Routing (RQ4)} --- from the intent routing pillar test matrix:
\begin{itemize}[leftmargin=*, nosep]
  \item \textbf{Mode classification accuracy:} per-mode accuracy (emotional support, practical education, policy navigation, mixed) before and after the full five-layer pipeline. Target: $\geq 90\%$ overall.
  \item \textbf{Confidence gating effect:} cases rerouted to keyword fallback; proportion corrected by fallback; net accuracy gain.
  \item \textbf{Safety override:} activation rate on hard-safety inputs (profanity, self-harm, crisis); false negative rate. Target: 0 false negatives.
  \item \textbf{Routing latency:} LLM router (P50, P95); keyword fallback path ($< 5$ ms expected); end-to-end routing decision latency.
\end{itemize}

\medskip
\textbf{[Table placeholder]} Routing evaluation summary: rows = LLM router pre-/post-gating, safety override, per-mode breakdown; columns = $n$ cases, correct, accuracy, latency P50/P95.
\end{tcolorbox}

\section{Expert Pilot Sessions}

\begin{tcolorbox}[colback=yellow!8!white, colframe=orange!70!black, title=\textbf{Placeholder --- Expert Pilot Sessions}, fonttitle=\bfseries, left=4pt, right=4pt, top=4pt, bottom=4pt]
\textit{Pending recruitment and session execution. This section will report findings from $n = 5$--8 domain experts (Spitex coordinators, geriatricians, home care nurses) evaluating system usability, response appropriateness, and ecological validity.}

\begin{itemize}[leftmargin=*, nosep]
  \item \textbf{Recruitment:} domain experts with direct informal caregiver support experience; purposive sampling across Spitex, geriatric care, and home care nursing roles.
  \item \textbf{Procedure:} structured walkthrough of 3--5 representative scenarios per coaching mode; think-aloud protocol; post-session semi-structured interview.
  \item \textbf{Measures:} System Usability Scale (SUS); appropriateness ratings per response (5-point Likert); open-ended feedback on personality modulation visibility and policy grounding accuracy.
  \item \textbf{Analysis:} descriptive statistics on SUS and appropriateness ratings; thematic analysis of qualitative feedback; inter-rater agreement on appropriateness ($\kappa$ target $\geq 0.75$).
\end{itemize}

\medskip
\textbf{[Table placeholder]} Expert pilot summary: SUS score, mean appropriateness per mode, key qualitative themes, inter-rater $\kappa$.
\end{tcolorbox}

\section{System Reliability (RQ5)}

\begin{tcolorbox}[colback=yellow!8!white, colframe=orange!70!black, title=\textbf{Placeholder --- System Reliability (RQ5)}, fonttitle=\bfseries, left=4pt, right=4pt, top=4pt, bottom=4pt]
\textit{Pending load test and audit-trail validation runs against the deployed Docker Compose stack. This section will report:}

\begin{itemize}[leftmargin=*, nosep]
  \item \textbf{Pillar test pass rate:} pass/fail across all 130+ test categories (personality detection, EMA stabilisation, intent routing, grounding verification, memory retrieval, audit write); failure-mode classification for any failures. Target: $\geq 95\%$.
  \item \textbf{Latency benchmarks:} end-to-end P50, P95, P99 under sequential load (1 user) and concurrent load ($\geq 10$ simultaneous sessions); LLM API call latency vs.\ pipeline overhead separately. Preliminary A/B results (Table~\ref{tab:perf}): parallel pipeline achieves ${\sim}8.1$\,s E2E (fresh) and ${\sim}3.5$\,s (cached) vs.\ ${\sim}9.5$\,s (sequential). Target: P50 $\leq 5$ s, P95 $\leq 8$ s.
  \item \textbf{JSON contract compliance:} proportion of inter-module messages passing Zod schema validation; failure types (missing fields, type violations, out-of-range OCEAN values); fallback activation rate on contract failure. Target: $\geq 99.5\%$.
  \item \textbf{Pipeline invariant adherence:}
    \begin{itemize}[nosep]
      \item Invariant 1 (universal personality detection): proportion of turns where detection ran.
      \item Invariant 2 (conditional retrieval): proportion of turns where retrieval was correctly triggered or suppressed by mode.
      \item Invariant 3 (mandatory grounding): proportion of policy responses with citations. All targets: 100\%.
    \end{itemize}
  \item \textbf{Audit trail completeness:} proportion of session turns with matching dual-write entries in session and audit log tables; gap or duplicate count.
  \item \textbf{Heuristic fallback frequency:} proportion of turns where API unavailable or timeout exceeded; correctness of heuristic OCEAN estimates on fallback turns.
\end{itemize}

\medskip
\textbf{[Table placeholder]} System reliability summary: rows = each metric; columns = result, target, status (pass/fail). To include latency percentiles, contract compliance rate, pillar test pass rate, invariant adherence counts, audit completeness.
\end{tcolorbox}

\chapter{Discussion}
\label{ch:discussion}

\section{Interpretation and Contributions}

\textbf{Completed finding (RQ1).} The personality detector achieves statistically significant Pearson correlations for all five OCEAN traits on the PERSONAGE benchmark ($n = 318$, mean $r = 0.47$, all $p < 10^{-12}$), with the style-aware ensemble reaching mean $r = 0.60$ on a balanced subset. The detector never inverts a trait's polarity, and the two-timescale EMA pipeline provides cross-session stability without requiring training data---a necessary prerequisite for the coaching quality claims that follow.

\textbf{Pending findings (RQ2--5).} The central empirical question---whether adding the personality-trait layer measurably improves coaching quality across emotional tone appropriateness, relevance and coherence, and emotional needs responsiveness---will be answered by the LLM-as-judge evaluation (RQ2). If the hypothesised effect sizes are confirmed (adaptive $\geq 4.0/5.0$ on ETA and REN; $\geq 20\%$ improvement over generic on REN), the memory-only versus full-system comparison will isolate the specific value of personality-trait modelling over the current state-of-the-art in LLM personalisation. Grounding safety (RQ3), routing accuracy (RQ4), and system reliability (RQ5) are evaluated through automated pillar test matrices and load benchmarks. The discussion below interprets the architecture and its design contributions; quantitative coaching quality claims are conditional on execution of these pending evaluations.

\textbf{Novelty against the literature.} The contributions can be positioned explicitly against prior work: \emph{Temporal stabilisation}---four-layer hierarchical EMA ($\alpha_s = 0.3$, $\alpha_p = 0.2$), confidence gating, prior-initialised session state; no prior system implements cross-session stabilisation without model training. \emph{Cross-session identity stability}---prior-initialised EMA anchors each session; identity drift \citep{Kovac2024} mitigated. \emph{Directive-based modulation}---interpretable, auditable, training-free; contrasts with RL/meta-learning \citep{denHengst2019, Lin2019, Shea2023}. \emph{Pipeline invariants}---three named, enforced invariants; no prior system. \emph{Hybrid intent routing}---five-layer router spanning task-oriented and open-domain; confidence gating, heuristic fallback, non-overridable safety. \emph{Content/style separation}---explicit architectural separation. \emph{Few-shot personality detection}---calibrated prompting; validated on PERSONAGE ($n = 318$, mean $r = 0.47$). \emph{Latency-optimised parallel pipeline}---concurrent detection and retrieval with detection caching; ${\sim}15\%$ E2E improvement on fresh requests and ${\sim}63\%$ on cached requests. \emph{Caregiver application}---first personality-aware coaching system for Swiss informal caregivers; directly addresses both dominant attrition drivers (emotional incongruence, lack of personalisation).

\textbf{Three personalisation layers as a coaching framework.} This work contributes a structured framework: transient state, content recall, and stable individual differences. Current systems implement only the first two. \textbf{Pipeline invariants as a design principle.} The three formal invariants represent a reusable design pattern for coaching systems that must be simultaneously adaptive and factually safe. \textbf{Content/style separation.} Big5Loop is, to our knowledge, the first system that must simultaneously deliver personality-modulated coaching and domain-grounded policy guidance---a tension that rule-based CBT chatbots sidestep by avoiding retrieval, and that RAG systems sidestep by omitting style modulation. \textbf{Contribution to caregiver support.} For Swiss caregivers, the system provides what no current digital tool offers: emotional resilience coaching adapted to individual trait profiles, practical education matched to learning preferences, and grounded Swiss policy navigation with citation-supported benefit guidance---integrated in a system that handles the common scenario where emotional distress and policy questions co-occur in the same turn.

\section{Reproducibility and Transferability}

A deliberate design objective is that the architecture can be independently reproduced and adapted to coaching domains beyond Swiss caregivers.

\textbf{Reproducibility} is enabled at four levels. At the \emph{component level}, all model parameters are fixed and versioned: Meta Llama 3 70B Instruct via NVIDIA API (\texttt{meta/llama3-70b-instruct}; temperature 0.1 for detection, temperature 0 for routing), in-session EMA $\alpha_s = 0.3$, cross-session EMA $\alpha_p = 0.2$, confidence gate $\tau_c = 0.4$, directive activation threshold $\theta_T = 0.15$. At the \emph{pipeline level}, Zod schema contracts between all modules (\texttt{@big5loop/contracts}) allow independent teams to reconstruct any module's expected inputs and outputs. At the \emph{evaluation level}, benchmark fixtures are versioned JSON, synthetic profiles are parameterised, and evaluation runs use fixed seeds. At the \emph{deployment level}, a single Docker Compose file with pinned image versions (N8N v2.12.0, PostgreSQL 16/pgvector, Node.js 20) reproduces the full stack. The dual audit trail enables full post-hoc reconstruction of any conversation---personality state, directives, routing decision, and grounding outcome at every turn.

\textbf{Transferability} is structural. The domain-specific components are: the policy corpus, the coaching scenario templates, and the coaching mode labels and descriptions. All other pipeline components---personality detection, EMA stabilisation, trait-to-directive mapping, hybrid intent routing, grounding verification, and audit infrastructure---contain no Swiss-specific assumptions. Replacing the policy corpus with a different knowledge base and updating scenario templates instantiates the architecture in a new coaching context without modifying the core pipeline.

\section{Limitations and Future Work}

\textbf{Limitations.} \emph{Pilot scope:} The current evaluation relies on automated test suites, benchmark infrastructure, and planned expert and simulated validation streams. Real caregiver validation ($n = 20$--30) is fully specified as a protocol but deferred due to recruitment feasibility. \emph{Language and dialect:} The prototype operates in English only; Swiss caregivers communicate primarily in Swiss German dialects, French, and Italian. Swiss German dialects are poorly supported by mainstream LLM backends trained on Standard German, creating a linguistic exclusion risk. \emph{Policy coverage:} The corpus covers selected Swiss social insurance domains; full coverage across all 26 cantons is needed for comprehensive guidance. \emph{Theoretical simplification:} The trait-to-directive mapping operationalises a simplified first-order version of the Zurich Model. \emph{LLM API latency:} Despite parallelisation and caching optimisations (\S\ref{sec:perf_opt}), the NVIDIA API remains the dominant latency contributor (${\sim}4$--5\,s per detection call, ${\sim}3$\,s per generation call); detection caching mitigates this for repeated messages but fresh requests remain API-bound. \emph{External API dependency:} The system depends on the NVIDIA API for inference and embeddings. \emph{Synthetic evaluation limitations:} Synthetic profiles cannot fully capture the variability and emotional complexity of real caregiving situations. \emph{LLM evaluator reliability:} LLM evaluators are susceptible to prompt sensitivity and can inherit biases; residual evaluator bias cannot be fully excluded \citep{Guo2023}.

\textbf{Future directions.} Real caregiver pilot ($n = 20$--30) with validated outcome measures; multilingual adaptation for Swiss German, French, and Italian; expanded policy corpus covering all 26 cantons; cross-domain transfer studies in non-caregiver coaching contexts; longitudinal evaluation over multiple weeks; empirical calibration of trait-to-directive mappings through iterative user studies; longitudinal attrition tracking---personality-aware modulation is hypothesised to reduce attrition specifically for high-Neuroticism users (survival analysis stratified by OCEAN profile); multimodal personality inference (e.g., voice prosody, typing dynamics) to reduce dependence on per-turn LLM detection; meta-RL policy upgrade path---once sufficient per-user interaction logs accumulate, the fixed-gain EMA could be replaced by a context-based meta-RL policy (e.g., PEARL-style task embeddings \citep{Rakelly2019}) or MAML-based personalisation \citep{Lin2019}, moving from rule-based directive mapping to learned policies without discarding the architectural invariants that ensure auditability.

\chapter{Conclusion}
\label{ch:conclusion}

Current digital coaching systems personalise through mood tracking and conversation memory but leave the stable personality traits that determine coaching receptivity entirely unaddressed---the third personalisation layer that human coaches apply implicitly. This study designed, implemented, and partially evaluated Big5Loop: a personality-aware coaching architecture that closes this gap by inferring Big Five traits from dialogue, stabilising them through EMA smoothing, and modulating coaching behaviour across four modes governed by three pipeline invariants ensuring universal personality detection, conditional retrieval, and mandatory grounding verification.

External validity of the personality detector was established against the PERSONAGE benchmark ($n = 318$, mean $r = 0.47$, all $p < 10^{-12}$; ensemble variant mean $r = 0.60$), confirming that calibrated few-shot prompting---not raw LLM capability---drives detection quality. Evaluations of coaching quality, grounding safety, routing accuracy, and system reliability (RQ2--5) are fully pre-specified and constitute the next phase of the research programme.

The architecture demonstrates that personality-aware modulation and citation-grounded domain knowledge can coexist in a single auditable pipeline. A partially parallelised orchestration layer (concurrent detection and retrieval) with detection caching achieves practical response latencies (${\sim}8$\,s fresh, ${\sim}3.5$\,s cached) despite dependence on external LLM APIs. Using Swiss informal caregivers as the pilot application, it provides the first implementation of personality-aware coaching for this population and the first formal treatment of personality modulation and factual grounding as simultaneously required but architecturally separated pipeline properties. The modular design---with explicit separation between coaching style adaptation and factual content delivery, reproducible contracts, and transferable pipeline components---is intended for independent reproduction and adaptation to other human-centred coaching domains.

% ========== REFERENCES ==========
% Per HSLU §5.2: APA 7 format, alphabetical order
\bibliographystyle{apalike}
\bibliography{manuscript_refs}

% ========== APPENDICES ==========
\appendix

% ── Appendix A: Technology Stack ─────────────────────────────────────────────
\chapter{Full Technology Stack}
\label{app:stack}

\begin{table}[H]
\centering\small
\caption{Big5Loop technology stack by layer}
\begin{tabularx}{\textwidth}{@{}llX@{}}
\toprule
Layer & Component & Version / Notes \\
\midrule
\textbf{Frontend} & Next.js & 14.2.0 (App Router, Turbopack dev server); Zustand state management; Tailwind CSS \\
\textbf{API Gateway} & Next.js API routes & TypeScript; Zod schema validation on all I/O contracts \\
\textbf{Auth} & JWT (HS256) & HTTP-only cookies; session isolation per user \\
\textbf{Orchestration} & N8N v2.12.0 (self-hosted) & JS task runners; parallel detection/retrieval pipeline; detection cache \\
\textbf{Language Model} & Meta Llama 3 70B Instruct & \texttt{meta/llama3-70b-instruct} via NVIDIA API; temperature 0.1 (detection), 0 (routing/evaluation) \\
\textbf{Embeddings} & NVIDIA NIM & \texttt{nv-embedqa-e5-v5}; 1024-dimensional vectors; used for pgvector retrieval \\
\textbf{Database} & PostgreSQL 16 & pgvector extension; 11 tables (see Appendix~\ref{app:schema}) \\
\textbf{Vector search} & pgvector & HNSW index on \texttt{policy\_chunks.embedding} and \texttt{personality\_memory\_embeddings.embedding} \\
\textbf{Full-text search} & PostgreSQL GIN & \texttt{to\_tsvector('english', content || ' ' || title)} index on \texttt{policy\_chunks} \\
\textbf{Contracts} & Zod schemas & Shared \texttt{@big5loop/contracts}; validated at API boundary and inter-module \\
\textbf{Deployment} & Docker Compose & Single-file stack: frontend, N8N, PostgreSQL; pinned image versions \\
\textbf{Reverse proxy} & Nginx & TLS termination; upstream routing to Next.js and N8N \\
\bottomrule
\end{tabularx}
\end{table}

% ── Appendix B: Trait-to-Directive Mapping ───────────────────────────────────
\chapter{Complete Trait-to-Directive Mapping}
\label{app:directives}

Directives activate when the in-session EMA estimate $T_{\text{sess},i}^{(n)}$ exceeds $\pm\theta_T = 0.15$. Multiple directives may be active simultaneously. When no trait exceeds the threshold, the default directive (warm, attentive, supportive, open-ended questions) applies.

\begin{table}[H]
\centering\small
\caption{Full trait-to-directive mapping with activation conditions and coaching intent}
\begin{tabularx}{\textwidth}{@{}llXl@{}}
\toprule
Trait & Direction & Behavioural directive and coaching intent & Activation \\
\midrule
Neuroticism (N) & High & Offer extra emotional comfort; explicitly acknowledge expressed anxieties before providing information; use validating language; avoid direct challenges & $T_N > 0.15$ \\
Neuroticism (N) & Low & Reinforce stability and competence; use confident, matter-of-fact tone; proceed to content without extended emotional check-in & $T_N < -0.15$ \\
Agreeableness (A) & High & Warm, empathetic, collaborative tone; emphasise shared goals and mutual support; use inclusive language (``we'', ``together'') & $T_A > 0.15$ \\
Agreeableness (A) & Low & Respectful but direct; reduce filler validation; deliver information efficiently without excessive warmth phrasing & $T_A < -0.15$ \\
Conscientiousness (C) & High & Structured, organised guidance; numbered steps; explicit timelines and action items; signal completeness (``here are all three steps'') & $T_C > 0.15$ \\
Conscientiousness (C) & Low & Reduced structure; conversational framing; highlight the single most important next action; avoid overwhelming checklists & $T_C < -0.15$ \\
Openness (O) & High & Invite exploration; surface multiple perspectives; use analogies and contextual framing; welcome follow-up questions & $T_O > 0.15$ \\
Openness (O) & Low & Concrete, practical guidance; avoid abstract framing; stick to established procedures; minimise optionality & $T_O < -0.15$ \\
Extraversion (E) & High & Energetic, sociable tone; positive reinforcement; acknowledge progress enthusiastically & $T_E > 0.15$ \\
Extraversion (E) & Low & Calm, reserved tone; avoid overly enthusiastic phrasing; give space for reflection; shorter responses preferred & $T_E < -0.15$ \\
\midrule
\multicolumn{3}{l}{\textit{Default (no trait active)}: warm, attentive, supportive; include one open-ended question} & all $|T_i| \leq 0.15$ \\
\bottomrule
\end{tabularx}
\end{table}

% ── Appendix C: Prompt Templates ─────────────────────────────────────────────
\chapter{Prompt Templates}
\label{app:prompts}

\begin{tcolorbox}[colback=gray!6!white, colframe=gray!50!black, title=\textbf{Placeholder --- Prompt Templates}, left=4pt, right=4pt]
\textit{To include the following four prompt templates extracted from the N8N workflow:}
\begin{itemize}[nosep]
  \item \textbf{Personality detection prompt}: system message defining trait-specific linguistic markers for each OCEAN dimension; seven few-shot examples (utterance $\to$ JSON scores); input slot for current user message.
  \item \textbf{Behavioural regulation prompt}: system message injecting active directive set $\mathcal{D}^{(n)}$ into response generation instructions; directive-to-style mappings per active trait.
  \item \textbf{Response generation prompt}: mode-specific system message (one per coaching mode); memory context injection; grounding instruction (policy mode only).
  \item \textbf{LLM intent router prompt}: system message and JSON output schema for mode classification (mode, confidence, retrieval\_flag, routing\_reason).
\end{itemize}
\end{tcolorbox}

% ── Appendix D: N8N Workflow Node Inventory ───────────────────────────────────
\chapter{N8N Workflow Node Inventory}
\label{app:workflow}

\begin{tcolorbox}[colback=gray!6!white, colframe=gray!50!black, title=\textbf{Placeholder --- N8N Workflow Node Inventory}, left=4pt, right=4pt]
\textit{To include: table listing all nodes in the V2 N8N workflow with node type, function, input/output contract, and key configuration parameters. Key V1 $\to$ V2 changes: (1) parallel branching from ``Merge Previous State'' into Detection and Retrieval paths with downstream merge node; (2) detection cache via \texttt{\$getWorkflowStaticData('global')} replacing \texttt{globalThis}; (3) reduced API timeouts (detection 12\,s, generation 20\,s) and \texttt{max\_tokens} (detection 50, generation 80--150); (4) GIN full-text index on \texttt{policy\_chunks(content || title)}. Estimated 20--25 nodes across the parallel pipeline.}
\end{tcolorbox}

% ── Appendix E: Database Schema ───────────────────────────────────────────────
\chapter{Database Schema}
\label{app:schema}

\begin{table}[H]
\centering\small
\caption{PostgreSQL database tables --- column specifications}
\begin{tabularx}{\textwidth}{@{}llX@{}}
\toprule
Table & Primary key & Key columns \\
\midrule
\texttt{users} & \texttt{id} UUID & \texttt{email}, \texttt{password\_hash}, \texttt{display\_name}, \texttt{locale}, \texttt{canton}, \texttt{created\_at} \\
\texttt{chat\_sessions} & \texttt{session\_id} UUID & \texttt{user\_id} (FK), \texttt{status}, \texttt{locale}, \texttt{canton}, \texttt{created\_at} \\
\texttt{user\_personality\_profiles} & \texttt{user\_id} UUID & \texttt{ocean\_scores} JSONB, \texttt{confidence} JSONB, \texttt{total\_turns}, \texttt{stable} BOOLEAN, \texttt{last\_updated} \\
\texttt{conversation\_turns} & (\texttt{session\_id}, \texttt{turn\_index}) & \texttt{user\_msg}, \texttt{assistant\_msg}, \texttt{mode}, \texttt{latency\_ms}, \texttt{created\_at} \\
\texttt{personality\_states} & (\texttt{session\_id}, \texttt{turn\_index}) & \texttt{ocean\_json} JSONB, \texttt{confidence\_json} JSONB, \texttt{stable} BOOLEAN, \texttt{ema\_alpha} \\
\texttt{policy\_evidence} & (\texttt{session\_id}, \texttt{turn\_index}, \texttt{source\_id}, \texttt{chunk\_id}) & \texttt{title}, \texttt{url}, \texttt{excerpt\_hash} \\
\texttt{performance\_metrics} & (\texttt{session\_id}, \texttt{turn\_index}, \texttt{stage}) & \texttt{status}, \texttt{duration\_ms}, \texttt{error\_code} \\
\texttt{personality\_memory\_embeddings} & \texttt{id} BIGSERIAL & \texttt{session\_id}, \texttt{turn\_index}, \texttt{memory\_text}, \texttt{embedding} VECTOR(1024) \\
\texttt{policy\_chunks} & \texttt{id} SERIAL & \texttt{source\_id}, \texttt{chunk\_id}, \texttt{title}, \texttt{content}, \texttt{url}, \texttt{embedding} VECTOR(1024), \texttt{metadata} JSONB \\
\texttt{audit\_log} & \texttt{id} BIGSERIAL & \texttt{request\_id}, \texttt{session\_id}, \texttt{turn\_index}, \texttt{mode}, \texttt{ocean\_json} JSONB, \texttt{directives} JSONB, \texttt{grounding\_score}, \texttt{citation\_count}, \texttt{verifier\_status}, \texttt{routing} JSONB, \texttt{turn\_latency\_ms} \\
\texttt{feedback} & \texttt{id} UUID & \texttt{session\_id}, \texttt{turn\_index}, \texttt{rating}, \texttt{comment}, \texttt{created\_at} \\
\bottomrule
\end{tabularx}
\end{table}

\noindent\textit{Indexes}: \texttt{idx\_sessions\_user} (user lookup); \texttt{idx\_turns\_session\_turn}, \texttt{idx\_personality\_session\_turn} (turn-level joins); \texttt{idx\_policy\_chunks\_content\_trgm} (GIN full-text); HNSW vector indexes on both embedding columns; \texttt{idx\_audit\_log\_session\_turn}, \texttt{idx\_audit\_log\_request}, \texttt{idx\_audit\_log\_created} (audit queries).

% ── Appendix F: Benchmark Fixture Format ─────────────────────────────────────
\chapter{Benchmark Fixture Format}
\label{app:fixtures}

\begin{tcolorbox}[colback=gray!6!white, colframe=gray!50!black, title=\textbf{Placeholder --- Benchmark Fixture Format}, left=4pt, right=4pt]
\textit{To include: JSON fixture schema (fields: \texttt{id}, \texttt{mode}, \texttt{input}, \texttt{personality\_profile}, \texttt{expected\_directives}, \texttt{expected\_mode}, \texttt{grounding\_required}, \texttt{pass\_criteria}); full case inventory (130+ cases by category); pass/fail criteria per test suite (personality detection, routing, grounding, audit).}
\end{tcolorbox}

% ── Appendix G: Evaluation Rubric ─────────────────────────────────────────────
\chapter{Evaluation Rubric and Scoring Guidelines}
\label{app:rubric}

\begin{tcolorbox}[colback=gray!6!white, colframe=gray!50!black, title=\textbf{Placeholder --- Evaluation Rubric}, left=4pt, right=4pt]
\textit{To include: 5-point Likert rubric for each of the three coaching quality dimensions (ETA, R\&C, REN) with anchor descriptions per scale point; LLM-as-judge prompt template and CoT chain structure; human expert inter-rater calibration procedure; disagreement resolution protocol.}
\end{tcolorbox}

% ── Appendix H: Swiss Policy Corpus ──────────────────────────────────────────
\chapter{Swiss Policy Corpus}
\label{app:corpus}

\begin{tcolorbox}[colback=gray!6!white, colframe=gray!50!black, title=\textbf{Placeholder --- Swiss Policy Corpus}, left=4pt, right=4pt]
\textit{To include: source family table (IV/Invalidenversicherung, AHV, Ergänzungsleistungen, Spitex, Hilflosenentschädigung, cantonal supplements); chunk count and token statistics per source; authority tier classification (federal vs.\ cantonal); URL freshness monitoring schedule and update policy.}
\end{tcolorbox}

% ========== DECLARATION OF SOLE AUTHORSHIP ==========
% Per HSLU Appendix A
\chapter*{Declaration of Sole Authorship}
\addcontentsline{toc}{chapter}{Declaration of Sole Authorship}

I, Jiahua Duojie, hereby certify that the attached work, \emph{Personality-Aware LLM-Based Chatbot for Informal Caregiver Coaching: Design, Implementation, and Pilot Evaluation}, is wholly and completely my own and that I have indicated all the sources (printed, electronic, personal, etc.) that I have consulted. Any sections quoted from these sources are clearly indicated in quotation marks or are otherwise so declared. I further attest that I have included acknowledgement of the name(s) of any person(s) consulted in the course of preparing this assignment.

\vspace{2cm}
Signed: \underline{\hspace{5cm}}

Date: \underline{\hspace{5cm}}

\end{document}
